\bookchapter{第六十七章：打磨與發佈——測試、除錯與上架 App Store}

\chapterintro{App 做出來了，然後呢？這一章回答“從能跑”到“能交付”的三件事：怎麼高效地除錯找 bug、怎麼寫測試防迴歸、怎麼把 App 送上 App Store 讓全世界下載。上架不只是儀式感——它是把一個作品變成產品的分水嶺。}

\booksection{一、除錯三板斧：print、斷點、預覽}

\begin{center}
\begin{tabularx}{\linewidth}{|X|X|X|}
\hline
\textbf{工具} & \textbf{幹什麼} & \textbf{怎麼用} \tabularnewline
\hline
print + 控制檯 & 看執行日誌 & 程式碼裡 print("值=\textbackslash{}(x)")，Xcode 底部控制檯看輸出 \tabularnewline
斷點除錯 & 讓程式停在某行，看變數值 & 行號旁點一下出藍標，執行到那裡暫停，逐步走 \tabularnewline
SwiftUI 預覽 + 偵錯程式 & 實時看介面與檢視層級 & \#Preview 畫布 + 執行時的 View Hierarchy 檢查器 \tabularnewline
\hline
\end{tabularx}
\end{center}

\begin{center}\small\color{gray}除錯三板斧\end{center}

\begin{tcbinfo}{小貼士：崩潰了怎麼讀？}
Swift 崩潰時控制檯會列印 “Fatal error” 加原因，Xcode 會停在出錯那一行並高亮。真機上的歷史崩潰可以在 Xcode Window Organizer Crashes 裡看（需付費帳號）。關鍵詞優先搜 “錯誤原因第一句 + swiftui”——精確的問題，才有精確的答案。

\end{tcbinfo}

\booksection{二、測試：給程式碼上保險}

手測靠不住——改了這個功能，另一個悄悄壞了。測試就是“後悔藥”兼“防彈衣”。iOS 用 XCTest 框架，分兩種：

\begin{center}
\begin{tabularx}{\linewidth}{|X|X|X|X|}
\hline
\textbf{層級} & \textbf{測什麼} & \textbf{Target} & \textbf{速度} \tabularnewline
\hline
單元測試 & 純邏輯（計算、校驗） & MyFirstAppTests & 秒級 \tabularnewline
UI 測試 & 真實介面行為 & MyFirstAppUITests & 分鐘級 \tabularnewline
手動測試 & 真人感受、體驗 & —— & 永遠需要 \tabularnewline
\hline
\end{tabularx}
\end{center}

\begin{center}\small\color{gray}iOS 測試層級（模板專案已預置兩個 Test Target）\end{center}

\begin{lstlisting}[language=Swift]
import XCTest

// 把"金額字串轉數字"這種邏輯抽成純函式
func parseAmount(_ input: String) -> Double? {
    guard let v = Double(input.trimmingCharacters(in: .whitespaces)),
          v > 0 else { return nil }
    return v
}

final class AmountTests: XCTestCase {
    func testEmptyInputReturnsNil() {
        XCTAssertNil(parseAmount(""))
    }
    func testNegativeReturnsNil() {
        XCTAssertNil(parseAmount("-5"))
    }
    func testNormalValueParses() {
        XCTAssertEqual(parseAmount("12.5"), 12.5)
    }
}
\end{lstlisting}

跑測試：CmdU。全綠之後再改程式碼，心裡有底。記帳本里的金額校驗、彙總計算，都該有單元測試罩著。

\booksection{三、簽名與打包：給 App 蓋上“印章”}

每個 iOS App 都必須簽名才能安裝——證明“這個包出自你，沒被篡改”。層級關係：

\begin{itemize}
  \item 學習階段：免費 Apple ID + Xcode 自動簽名（Personal Team），真機可跑，7 天續期一次；
  \item 分發階段：付費開發者帳號（99 美元/年）簽發證書與描述檔案，才能上架和長期安裝。
\end{itemize}

打包：Xcode 選單 Product Archive（歸檔），構建完成後在 Organizer 裡 Distribute App 選 App Store Connect 上傳。歸檔必須用 Generic iOS Device 作為目標（不是模擬器），這是新手常忘的一步。

\textbf{避坑提醒：證書丟了 = 更新斷檔}\par

簽名證書和私鑰一旦丟失，你無法再給舊 App 發更新——使用者只能裝“全新”應用。妥善備份證書（鑰匙串匯出 .p12）並記好密碼。好在蘋果提供了“自動簽名”和雲證書管理（Xcode 裡勾選 Automatically manage signing），能省掉大部分手工操作的坑。

\booksection{四、上架 App Store 全流程}

註冊開發者帳號

準備商店素材

上傳構建包

填隱私與稽核資訊

蘋果稽核（1-7 天）

上架

\begin{center}\small\color{gray}圖 5：App Store 上架六步走\end{center}

\begin{itemize}
  \item 註冊：developer.apple.com，個人帳號 99 美元/年；
  \item 素材：圖示（1024×1024，別帶圓角——蘋果自己切）、截圖（iPhone 各尺寸）、描述、隱私政策連結；
  \item 隱私問卷：如實申報收集了哪些資料——不聯網的記帳本可以"不收集任何資料"；
  \item 稽核：首次提交常有 1-7 天等待，被拒就按"稽核備註"逐條改——蘋果的拒絕理由通常寫得非常具體，照做即可；
  \item TestFlight：上架前先把測試版發給朋友內測（最多 10000 人），收集一輪反饋再正式提交，專業範兒十足。
\end{itemize}

\booksection{五、商店素材怎麼做：第一印象就是下載量}

商店頁是使用者“逛店”時看到的櫥窗。三個素材決定第一印象：

\begin{itemize}
  \item 圖示：1024×1024，別用截圖湊合。沒設計能力就用 AI 工具生成一個簡潔圖形，在桌面上縮小後還能一眼認出才算合格；
  \item 截圖：至少 2 套（iPhone 主流尺寸），展示最核心的功能場景；每張配一句人話文案（"一秒記一筆"比"高效的記帳解決方案"好十倍）；
  \item 描述：第一句說清"這是什麼 + 給誰用"，別堆砌"最專業、最強"——真誠比誇張有效。
\end{itemize}

\begin{tcbpractice}{來動手：把記帳本"交付"一次}
\booksection{六、崩潰上報與效能最佳化小清單}

App 上架前，這兩件事值得各花半小時：

\booksubsection{崩潰監控}

你測試時沒崩，不代表使用者手機不崩。蘋果在 Xcode Organizer 裡提供免費的崩潰報告（付費帳號可看真機崩潰日誌）；更省心的方案是接入第三方 SDK（如 Firebase Crashlytics）。沒有崩潰監控的上架，等於閉著眼睛開車。

\booksubsection{效能自查四條}

\begin{enumerate}
  \item 列表用 LazyVStack / List：長列表別用普通 VStack 硬排——懶載入是效能第一課；
  \item 圖片別硬塞原圖：用 AsyncImage（遠端圖）和系統快取，5MB 原圖足以卡爆老機型；
  \item 別在主執行緒乾重活：網路、大檔案都走 async/await；
  \item 用 Instruments 量一量：Xcode 自帶的效能分析器（Time Profiler、記憶體洩漏檢查）——先測量，再最佳化，別憑感覺。
\end{enumerate}

(1) 給 parseAmount 等邏輯補 3 個單元測試（CmdU 跑綠）；(2) 用斷點親手抓一個自己製造的崩潰並修好；(3) Product Archive 打一個包，走 TestFlight 邀請朋友內測；(4) 有餘力就註冊開發者帳號，把記帳本真正提交上架。前三步今晚就能做完。

\end{tcbpractice}

\begin{tcbsummary}{本章小結}
\begin{itemize}
  \item 除錯三板斧：print 日誌、斷點、SwiftUI 預覽與檢視檢查器
  \item XCTest：單元測試（CmdU）+ UI 測試 + 手動測試，模板已預置
  \item 簽名：免費帳號自動簽名（7 天續期）；付費帳號才能分發
  \item 打包：Generic iOS Device + Archive + Distribute；證書務必備份
  \item 上架六步：註冊 素材 上傳 隱私問卷 稽核 上架；TestFlight 先內測
\end{itemize}

\end{tcbsummary}
